%% ============================================================
%% Standalone interim draft for the supervisor (29 July 2026).
%% Contents: full TOC skeleton, Theory chapter, Method interim
%% summary, coding-table excerpt (S01–S34), List of Aids.
%% Preamble mirrors main.tex. main.tex stays the master file.
%% ============================================================
\documentclass[12pt, a4paper]{report}

\usepackage[utf8]{inputenc}
\usepackage[T1]{fontenc}
\usepackage[english]{babel}
\usepackage{csquotes}

\usepackage[style=apa, backend=biber]{biblatex}
\addbibresource{bib.bib}

\usepackage[a4paper, margin=2.5cm]{geometry}
\usepackage{setspace}
\onehalfspacing

\usepackage[protrusion=true,expansion=false]{microtype}

\usepackage{booktabs}
\usepackage{tabularx}
\usepackage{graphicx}
\usepackage{float}
\usepackage{longtable}
\usepackage{pdflscape}

\usepackage[hidelinks]{hyperref}

%% ============================================================
\begin{document}

% The citation-verification footnotes (verbatim source passages, see CLAUDE.md
% convention) are an internal checking device — suppressed in this draft only;
% sections/background.tex keeps them. Must come AFTER \begin{document}:
% hyperref/biblatex re-define \footnote at begin-document and would undo it.
\renewcommand{\footnote}[2][]{}

\begin{titlepage}
  \centering
  \vspace*{2cm}
  {\LARGE\bfseries Organizational Conditions for AI-Driven Firm Performance\par}
  \vspace{0.5cm}
  {\large A Systematic Review of Empirical Evidence, 2015--2026\par}
  \vspace{2cm}
  {\large Bachelor's Thesis --- interim draft for supervisor review\par}
  {\large Supervisor: PD Dr.\ Viktor Fredrich\par}
  \vspace{1cm}
  {\large Vienna University of Economics and Business (WU Wien)\par}
  \vspace{1cm}
  {\large Arthur Wienerroither\par}
  {\large Student ID: h12322363\par}
  \vspace{0.5cm}
  {\large 29 July 2026\par}
\end{titlepage}

\tableofcontents
\newpage

\chapter{Introduction}
\label{sec:intro}
\noindent\emph{This chapter follows. Research question and framing as approved in the proposal.}

\chapter{Theoretical Background}
\label{sec:background}

% Key concepts, definitions, and theoretical frameworks.
% Review of relevant literature; establish the state of research.

\section{Competitive advantage and firm performance}
\label{sec:constructs}

Competitive advantage is measured in 15 of the 67 studies reviewed here, and the
remaining 52 report performance outcomes only. Barney separates the two constructs by
what rivals are able to do: a firm holds a competitive advantage while it runs a
value-creating strategy no rival, actual or potential, is running at the same time, and
that advantage counts as sustained only once those rivals also cannot reproduce what the
strategy yields \parencite[p.~102]{barney1991firm}\footnote{[Barney (1991), p.~102] ``In
this article, a firm is said to have a \emph{competitive advantage} when it is
implementing a value creating strategy not simultaneously being implemented by any current
or potential competitors. A firm is said to have a \emph{sustained competitive advantage}
when it is implementing a value creating strategy not simultaneously being implemented by
any current or potential competitors \emph{and} when these other firms are unable to
duplicate the benefits of this strategy.''}. Sustained in that sense says nothing about
elapsed time. It asks whether the advantage is still there after rivals have given up
trying to copy it \parencite[p.~102]{barney1991firm}\footnote{[Barney (1991), p.~102]
``Rather, whether or not a competitive advantage is sustained depends upon the possibility
of competitive duplication. Following Lippman and Rumelt (1982) and Rumelt (1984), a
competitive advantage is sustained only if it continues to exist after efforts to
duplicate that advantage have ceased.''}. This review therefore keeps the two apart
throughout, because a firm can raise its margins without holding anything a rival is
unable to copy.

Where both constructs appear, the separation is visible in the survey items.
Chatterjee et al. ask about rivals directly: one of their competitive-advantage items
has respondents affirm that successful implementation of AI-based customer relationship
management ``will help an organization to win over its competitor which has not yet
implemented AI-CRM for B2B relationship management'' \parencite[p.~215]{chatterjee2021effect}\footnote{[Chatterjee et al. (2021),
p.~215] ``I think that successful implementation of AI-CRM will help an organization
to win over its competitor which has not yet implemented AI-CRM for B2B relationship
management.''}. Their performance items stay inside the firm and ask whether the same
system improves ``operational efficiency'' and speeds up
decisions \parencite[p.~215]{chatterjee2021effect}\footnote{[Chatterjee et al. (2021),
p.~215] ``I believe that successful implementation of AI-CRM will help the organization
to improve its operational efficiency.'' / ``I think AI-CRM helps in quick decision
making which helps the organization to improve its operational performance.''}.

Mehta et al. give competitive advantage a position between adoption and results: they
test it as the channel through which AI adoption reaches performance, and predict a
positive effect on performance in its own
right \parencite[p.~550]{mehta2025driven}\footnote{[Mehta et al. (2025), p.~550]
``H8. Competitive advantage has a significant and positive impact upon firm
performance.'' / ``H9. Competitive advantage has a mediating influence on adoption of
AI technologies and firm performance.''}.

Not every study that names competitive advantage builds its own instrument. Hossain
et al. measure sustained competitive advantage with items carried over from earlier
work \parencite[p.~247]{hossain2022marketing}\footnote{[Hossain et al. (2022), p.~247]
``sustained competitive advantage items were taken from Cao et al. (2019).''}. The
construct label alone says little about what a study measured; the items decide.

\section{Artificial intelligence as a general-purpose technology}
\label{sec:gpt}

Economists file AI under general-purpose technologies (GPTs), and that classification
already predicts an uneven payoff. A GPT earns the label because it makes ``new and
complementary production methods'' available, which can raise output over
time \parencite[p.~188]{czarnitzki2023artificial}\footnote{[Czarnitzki et al. (2023),
p.~188] ``The main feature of general-purpose technology is to enable new and
complementary production methods that may increase productivity over time (Bresnahan
and Trajtenberg 1995; Bresnahan et al., 2002; Brynjolfsson and Hitt 2003; Cardona
et al., 2013).''}. Those complements are the expensive part. Before output moves, a
firm has to redesign how work is organised, retrain staff, rework its software and
accumulate managerial experience, and because none of that accumulation is booked as
an asset, measured productivity dips before it
climbs \parencite[pp.~333--334]{brynjolfsson2021productivity}\footnote{[Brynjolfsson
et al. (2021), pp.~333--334] ``realizing their potential also requires large intangible
investments and a fundamental rethinking of the organization of production itself.
Firms must create new business processes, develop managerial experience, train workers,
patch software, and build other intangibles. This raises productivity measurement
issues because intangible investments are not readily tallied on a balance sheet or in
the national accounts.'' / [p.~334] ``The error in measured total factor productivity
growth therefore follows a J-curve shape, initially dipping while the investment rate in
unmeasured capital is larger than the investment rate in other types of capital, then
rising as growing intangible stocks begin to contribute to measured
production.''}. The lag is years, not quarters.

Two studies in this corpus show how differently that plays out at firm level. On German
innovation-survey panel data, Czarnitzki et al. measure a clear productivity gain from
AI use, for the yes/no measure and for the intensity measure
alike \parencite[pp.~189, 201]{czarnitzki2023artificial}\footnote{[Czarnitzki et al.
(2023), p.~189] ``We analyze cross-section as well as panel data from the German part of
the European Commission's Community Innovation Survey (CIS).'' / [p.~201] ``We found
that employing AI technologies has a positive and significant impact
on firm productivity. In particular, we showed that both the use of AI and the intensity
with which firms exploit the potential of AI significantly increase both sales and
value-added.''}. Fontanelli et al. ask a different question of French firms and find
that AI use raises the \emph{volatility} of productivity growth; the effect sits with
firms that buy AI from outside providers, and it weakens where more of the payroll is
made up of ICT engineers and
technicians \parencite[pp.~1--2]{fontanelli2025predictive}\footnote{[Fontanelli et al.
(2025), p.~1] ``This paper shifts the focus and investigates the impact of predictive AI
on the volatility of firms' productivity growth rates.'' / ``Using firm-level data from
the 2019 French ICT survey, we provide robust evidence that AI use is associated with
increased volatility.'' / [pp.~1--2] ``Our results show that heightened volatility is concentrated among AI
buyers, whereas firms that develop AI internally experience no such association.
Finally, we find that the AI-volatility link among `AI buyers' is mitigated in firms
with a higher share of ICT engineers and technicians, suggesting that AI's successful
integration requires complementary human capital.''}. Across the two results, the
reading this review takes forward is that the same AI investment can stand behind
either outcome; what differs between the firms is the complements.

\section{The resource-based view and the imitability problem}
\label{sec:rbv}

Across the studies reviewed here, the resource-based view (RBV) is the recurring reason
given for expecting AI to pay off. Barney sets four conditions on a resource before it
can hold an advantage in place: it has to be valuable, rare among competitors and
imperfectly imitable, and no substitute may do the same strategic
job \parencite[pp.~105--106]{barney1991firm}\footnote{[Barney (1991), pp.~105--106] ``To
have this potential, a firm resource must have four attributes: (a) it must be valuable,
in the sense that it exploit opportunities and/or neutralizes threats in a firm's
environment, (b) it must be rare among a firm's current and potential competition, (c) it
must be imperfectly imitable, and (d) there cannot be strategically equivalent substitutes
for this resource that are valuable but neither rare or imperfectly imitable.''
Typographical errors as in the original.}. Alwakid and Dahri restate the four conditions
as the VRIN
acronym \parencite[p.~14]{alwakid2025harnessing}\footnote{[Alwakid \& Dahri (2025), p.~14]
``Traditionally, RBV posits that firms achieve sustained competitive advantage by
strategically deploying valuable, rare, inimitable, and non-substitutable (VRIN)
resources (Barney, 1991; Wang \& Zhang, 2025).''}. That is a shorthand Barney himself comes close
to when he lists ``rareness, inimitability, non-substitutability'' among the
characteristics that can qualify a firm attribute as a source of advantage, adding that
the attribute becomes a resource only once it exploits an opportunity or neutralizes a
threat \parencite[p.~106]{barney1991firm}\footnote{[Barney (1991), p.~106] ``Firm
attributes may have the other characteristics that could qualify them as sources of
competitive advantage (e.g., rareness, inimitability, non-substitutability), but these
attributes only become resources when they exploit opportunities or neutralize threats in
a firm's environment.''}. Hossain et al. put the same argument
in terms of what a firm does with what it owns, and tie the long-run case to holdings
rivals cannot match \parencite[pp.~241--242]{hossain2022marketing}\footnote{[Hossain
et al. (2022), pp.~241--242] ``The resource-based view (RBV) portrays how a firm can
optimally use its resources, which are non-substitutable, unique and valuable for
achieving sustained competitive advantage (Wernerfelt, 1984).''}.

Purchased AI fails that test, and the RBV literature says so plainly. Krakowski et al.
note that AI costs almost nothing to reproduce and that little stands in the way of
copying it, so where AI substitutes a human capability, the RBV predicts that the
advantage resting on that capability will
erode \parencite[p.~1426]{krakowski2023artificial}\footnote{[Krakowski et al.
(2023), p.~1426] ``When AI substitutes humans' cognitive capabilities, the RBV expects
the advantage that these capabilities provided to erode (Peteraf \& Bergen, 2003). This
is because, as a technology resource, AI has close to zero marginal reproduction costs
and few imitation barriers (Brynjolfsson \& McAfee, 2014, p.~31).''}. Their own evidence
relocates the advantage rather than denying it: what resists imitation is the combination
a firm builds, not either side of
it \parencite[p.~1446]{krakowski2023artificial}\footnote{[Krakowski et al. (2023),
p.~1446] ``Our findings show that humans remain sources of sustainable advantage by
combining and integrating their capabilities with machine capabilities into uniquely
complementary bundles (Argyres \& Zenger, 2012).'' / ``the locus of sustainable advantage
is neither human nor machine capabilities, but what humans do with these capabilities.''}.

Two consequences follow for this review. The RBV speaks to competitive advantage and not
to performance, so a study reporting higher margins has not thereby shown that anything
is defensible against rivals; the separation of the two constructs in
Section~\ref{sec:constructs} is a theoretical requirement, not only a coding convention.
And if the resource itself is imitable, the variation worth studying lies in what
surrounds it, which is why the conditions under which AI pays off are treated here as the
object of analysis rather than as controls.
% Barney (1991) wird ueber literature/barney1991firm.txt zitiert (verifiziertes
% Text-Sidecar zum Bildscan, Seitenmarker = Journalseiten). Alle Passagen am
% Seitenbild gegengelesen: S. 102 (Definitionen, in Abschnitt 2.1) sowie
% S. 105-106 und S. 106 (vier Attribute, hier).

\section{Automation and augmentation}
\label{sec:autoaug}

The same investment can be deployed in two ways, and the strategy literature keeps them
apart. Raisch and Krakowski define automation as machines taking a task over from a
human, and augmentation as humans and machines working on a task
together \parencite[p.~192]{raisch2021artificial}\footnote{[Raisch \& Krakowski (2021),
p.~192] ``Whereas automation implies that machines take over a human task, augmentation
means that humans collaborate closely with machines to perform a task.''}. Their argument
against the popular advice to prefer augmentation is that the two cannot be held apart in
practice. For any single task at any single moment the firm picks one approach, and that
choice makes the two contradictory; widen the frame to several tasks over several years
and each turns out to produce the other, which makes them
interdependent \parencite[pp.~195--197]{raisch2021artificial}\footnote{[Raisch \&
Krakowski (2021), p.~195] ``Automation and augmentation are contradictory, because
organizations choose either one or the other approach to address a given task at a
specific point in time.'' / ``If we increase our analysis's temporal scale (from one
point in time to its evolution over time) and spatial scale (from one to multiple
tasks), we comprehend that, in the management domain, the two AI applications are not
only contradictory but also interdependent.'' / [p.~196] ``Consequently, augmentation
may enable a transition to automation over time.'' / [p.~197] ``Automation in one task
`spills over,' enabling adjacent tasks' augmentation.''}.

Corpus studies pick the distinction up as a design choice with different returns. Zebec
and Indihar {\v S}temberger treat the two as the use cases through which AI adoption
reaches performance, automation paying off in efficiency and augmentation in what people
and machines manage
jointly \parencite[p.~7]{zebec2024creating}\footnote{[Zebec \& Indihar {\v S}temberger
(2024), p.~7] ``Automation involves machines taking over tasks from humans, leading to
performance gains via resource, cost and information-processing efficiencies. In
contrast, augmentation involves humans collaborating with machines, enhancing both human
and machine skills and productivity.''}. Their model then returns no significant direct
effect of process automation on process performance once tasks become complex, and the
recommendation they draw from it is a boundary condition: automate what is routine,
augment what is
ambiguous \parencite[pp.~16--17]{zebec2024creating}\footnote{[Zebec \& Indihar
{\v S}temberger (2024), pp.~16--17] ``The results show that BPA has no significant direct
effect on BPP for more complex tasks in terms of efficiency of execution or scalability.''
/ ``In summary, automate routine, well-structured tasks and augment complex, ambiguous
ones.''}. Where the line falls is itself unstable. In the fragrance industry the machine
took over part of the perfumers' problem solving and could not touch the part that
depends on smelling a scent and judging how people will react to
it \parencite[p.~1445]{krakowski2023artificial}\footnote{[Krakowski et al. (2023),
p.~1445] ``While machines substitute perfumers regarding certain problem solving
activities, they are unable to match the perfumers' unique ability to smell fragrances
and predict the human emotions they trigger. This example illustrates that AI is likely
to substitute some, but not all of complex business tasks' activities.''}.

\section{Situated AI}
\label{sec:situated}

Kemp starts from the puzzle the two previous sections produce: AI is available to
everyone, so buying it should not yield an advantage, yet firms plainly differ in what
they get out of it. His answer is situated AI, which he defines as AI hemmed in by the
firm's own experience, structure and
relationships \parencite[p.~618]{kemp2024competitive}\footnote{[Kemp (2024), p.~618]
``The present paper begins to resolve this puzzle by developing a theory of situated
AI---AI whose agency is circumscribed in a firm's experiential, structural, and
relational systems.''}. Three properties of the technology stand in the way, and the
first of them is the imitability problem from Section~\ref{sec:rbv} under another name:
AI is generic, explicit and myopic. Being generic puts AI on a level with general human
capital, which competitors can acquire and put to work much as the adopter
does \parencite[p.~620]{kemp2024competitive}\footnote{[Kemp (2024), p.~620] ``The generic,
explicit, and myopic nature of AI all endanger its firm-level benefits, but are uniquely
difficult to surmount because AI is both a machine and agentic (Smith \& Lewis, 2011).
For example, AI's generic nature is akin to general human capital. General human capital
cannot typically underly interfirm advantages because a firm's competitors can usually
acquire and deploy that knowledge in ways that closely mimic corresponding actions of the
focal firm.''}.

What a firm does control is the situating. Kemp names three activities: grounding, which
picks the experiences a firm's own AI learns from; bounding, which works on the
experiences feeding a competitor's AI; and recasting, which keeps reworking the
algorithms and the routines around them so the AI stays aligned with the rest of the
firm's activities \parencite[pp.~618--619]{kemp2024competitive}\footnote{[Kemp (2024),
pp.~618--619] ``Grounding involves orchestrating which experiences one's AI will be
allowed to learn from across the organization. Bounding involves efforts to shape the
experiences anchoring a competitor's AI. Recasting involves orchestrating the continual
adaptation of algorithms and their surrounding routines to enhance AI's alignment with
interdependent activities in a firm.''}. The vocabulary is conditional throughout. None
of the three is a property of the technology, and each is something a firm may or may not
do.

Empirical work in the corpus starts from the same reading. Shi et al. take the mixed
results of earlier work as a sign that what AI yields depends on intangible complements,
on governance, and on managerial capability bound to its context, and they test that
claim on Chinese listed firms, with data assets as the channel and managerial capability
as the moderator \parencite[p.~1]{shi2025value}\footnote{[Shi et al. (2025), p.~1] ``Drawing on a
comprehensive dataset of Chinese listed companies spanning 2010--2022, this study
investigates how AI adoption influences enterprise value through the mediating role of
data assets and the moderating effect of managerial capability.'' / ``These mixed
outcomes suggest that the value-creating capacity of AI depends on more than technological
sophistication alone; it hinges on complementary intangible resources, effective
governance structures, and context-specific managerial capabilities---factors that current
empirical research has only partially explored.''}. Kemp's paper is conceptual, so what it
offers are propositions rather than findings. They are the propositions this review is
positioned to examine, because the conditions Kemp derives are the conditions the coded
studies report.


% Interim short version of the Method chapter for the supervisor draft of
% 29 July 2026 (compiled via BA_Zwischenstand_2026-07-29.tex, NOT via main.tex).
% The full chapter will be written in sections/method.tex; every number here is
% frozen method documentation (see CLAUDE.md), not derived from the coding table.
\chapter{Method}
\label{sec:method-summary}

\noindent\emph{Interim summary. The full chapter --- search documentation, screening protocol, coding procedure, and PRISMA flow diagram --- follows in the next version of this draft.}

\bigskip

\noindent The review follows the protocol announced in the proposal. The documented Scopus query was run on 17 July 2026 and returned 432 unique records: peer-reviewed journal articles in English, published 2015--2026, in the Scopus subject areas business, economics, and decision sciences. The query finds all three benchmark studies defined in advance (Babina et al., 2024; Krakowski et al., 2023; Lui et al., 2022).

For journal quality, the strict AJG~4*/4 threshold of the proposal was replaced by a broader, uniform inclusion rule: a record stays in when its journal is listed in the AJG~2024 with a rating of at least~2. One fixed ranking edition applies to all records, so the rule is transparent and non-selective. This filter reduced the corpus from 432 to 174 records; most of the excluded records appeared in journals not listed in the ranking.

Title and abstract screening tested four criteria: the study is empirical, works at the firm level, treats firms' AI investment or adoption as its object, and measures a performance or competitive-advantage outcome. This step reduced the 174 records to 81. Reading the full texts removed another fourteen studies, twelve on content grounds and two that could not be retrieved. The final corpus holds 67 studies. Every screening decision is recorded with a reason code, and the full chain (432 $\rightarrow$ 174 $\rightarrow$ 81 $\rightarrow$ 67) will be reported as a PRISMA-style flow diagram.

Each included study was coded into the state-of-the-art table announced in the proposal; an excerpt is attached in the Appendix (Tables~\ref{tab:coding-a1} and~\ref{tab:coding-a2}). Following the supervisor's feedback, the scheme separates superior firm performance from competitive advantage as distinct outcome constructs, each with its own measurement column.


\chapter{Results}
\label{sec:results}
\noindent\emph{In progress. Part one, the descriptive mapping of the corpus, is drafted; part two, the synthesis of conditions, is the next writing step.}

\chapter{Discussion}
\label{sec:discussion}
\noindent\emph{This chapter follows.}

\chapter{Conclusion and Future Outlook}
\label{sec:conclusion}
\noindent\emph{This chapter follows.}

\newpage
\printbibliography[heading=bibintoc, title={References}]

\newpage
\appendix
\chapter*{Appendix}
\addcontentsline{toc}{chapter}{Appendix}

\section*{Coding table (excerpt)}
\label{app:coding}

\noindent An excerpt of the coding table: the 25 studies published in the highest-ranked journals (AJG 4* and 4, then 3; within a rating ordered by year) plus, at the end, the first three AJG~2 studies as examples of the lowest included rating tier --- 28 of 67 coded studies in total. Split into study characteristics (Table~\ref{tab:coding-a1}) and findings (Table~\ref{tab:coding-a2}); study IDs refer to the full table.

% AUTO-GENERATED by corpus/make_appendix_table.py — do not edit by hand.
\begingroup
\renewcommand{\thetable}{A.\arabic{table}}
\setcounter{table}{0}
\begin{landscape}
\scriptsize
\begin{longtable}{@{}l p{2.6cm} p{3.4cm} p{1.8cm} p{5.0cm} p{2.1cm} p{6.2cm}@{}}
\caption{Included studies: characteristics and design (excerpt: 28 of 67 studies --- the 25 in the highest-ranked journals (AJG 4*, 4, 3) plus the first three AJG 2 studies; within a rating by year). Source: author's compilation.}\label{tab:coding-a1}\\
\toprule
ID & Study & Journal (AJG) & Country & Sample & Method & AI investment measure \\
\midrule
\endfirsthead
\multicolumn{7}{@{}l}{\tablename~\thetable{} (continued)}\\
\toprule
ID & Study & Journal (AJG) & Country & Sample & Method & AI investment measure \\
\midrule
\endhead
\bottomrule
\endfoot
S09 & Mishra et al. (2022) & Journal of the Academy of Marketing Science (4*) & USA & 19,000 firm-year observations, US-listed firms (COMPUSTAT + EDGAR 10-Ks), simultaneous equations; PSM robustness (1,801 treatment / 2,251 control firms) & panel econometrics & 10-K text (share of AI keywords = AI focus) \\
\addlinespace[2pt]
S14 & Krakowski et al. (2023) & Strategic Management Journal (4*) & multi-country & same chess players across conventional, centaur and engine tournaments (controlled competitive setting) & panel econometrics & AI (engine) adoption in competitive interaction \\
\addlinespace[2pt]
S17 & Babina et al. (2024) & Journal of Financial Economics (4*) & USA & 1,993 US Compustat firms; resume-based AI workforce measure (Cognism), robust to job-postings measure; IV = university AI graduate supply & panel econometrics & resume-based (share of AI-skilled workers as firm AI investment) \\
\addlinespace[2pt]
S07 & Leoni et al. (2022) & International Journal of Operations and Production Management (4) & Italy & 120 senior executives of manufacturing firms, PLS-SEM & survey-SEM & survey construct (AI adoption/use) \\
\addlinespace[2pt]
S45 & Adams et al. (2026) & Journal of Monetary Economics (4) & USA & US firms matched to Lightcast job postings 2010Q1-2024Q1 (40,000+ job boards); AI-pricing jobs identified & panel econometrics & job postings (AI-skill + pricing keyword = AI pricing adoption) \\
\addlinespace[2pt]
S46 & Agag et al. (2026) & Tourism Management (4) & UK & 1,206 UK listed tourism/travel/hospitality firms, 7,539 firm-years 2018-2024; two-step system GMM & panel econometrics & AI human capital (AI-skilled workforce measure) \\
\addlinespace[2pt]
S52 & Filieri et al. (2026) & Journal of Product Innovation Management (4) & China & Chinese listed manufacturers + WTO tariff records + customs data 2015-2022; IV analysis & panel econometrics & archival AI integration measure (tariff-induced) \\
\addlinespace[2pt]
S02 & Baabdullah et al. (2021) & Industrial Marketing Management (3) & Saudi Arabia & 392 B2B SMEs, cross-section, CB-SEM & survey-SEM & survey construct (acceptance of AI practices) \\
\addlinespace[2pt]
S03 & Chatterjee et al. (2021) & Industrial Marketing Management (3) & India & 349 managers from 16 BSE-listed firms (11 manufacturing, 5 service), response rate 51.6\%, PLS-SEM & survey-SEM & survey construct (AI-CRM implementation) \\
\addlinespace[2pt]
S04 & Chatterjee et al. (2022) & Journal of Business Research (3) & India & 312 usable manager responses from 14 BSE-listed firms (10 large/4 SME; telecom, IT, automobile, retail, pharma), Jan-Feb 2020, RR 47.3\%, PLS-SEM & survey-SEM & survey construct (adoption of AI-embedded CRM system) \\
\addlinespace[2pt]
S05 & Hossain et al. (2022) & Industrial Marketing Management (3) & Bangladesh & 257 usable manager responses (1,953 approached, 278 qualified), RMG export manufacturers, research-firm panel/Qualtrics, random sampling; time-lagged robustness subsample n=93 & survey-SEM & survey construct (AI adoption on top of marketing analytics platform) \\
\addlinespace[2pt]
S06 & Lee et al. (2022) & Technovation (3) & South Korea & 160 high-tech ventures (from 1,248 ministry list, 300 responses), survey + audited financials & panel econometrics & survey (AI-adoption intensity) linked to financial records \\
\addlinespace[2pt]
S08 & Lui et al. (2022) & Annals of Operations Research (3) & multi-market & 119 AI investment announcements by 62 listed firms, Factiva search Jan 2015 - Feb 2019 & event study & announcements (AI investment announcements) \\
\addlinespace[2pt]
S12 & Czarnitzki et al. (2023) & Journal of Economic Behavior and Organization (3) & Germany & 5,851 firms (409 AI users, \textasciitilde{}7\%), German CIS 2018 (ZEW); cross-section OLS/2SLS-IV + panel subset with IV-FE, entropy balancing & panel econometrics & survey (CIS AI adoption dummy + intensity of AI use in business processes) \\
\addlinespace[2pt]
S15 & Samadhiya et al. (2023) & Industrial Marketing Management (3) & India & in-depth interviews with senior B2B managers + PLS-SEM on 142 responses & mixed & survey construct (AI-based partner relationship management implementation) \\
\addlinespace[2pt]
S16 & Wu \& Monfort (2023) & Psychology and Marketing (3) & Taiwan & 278 food firms (CEO/marketing managers), Taiwan, survey Mar-Aug 2020, RR 33\%; SEM + supplementary fsQCA & survey-SEM & survey construct (implementation of AI marketing strategy) \\
\addlinespace[2pt]
S18 & Cannas et al. (2024) & International Journal of Production Research (3) & Italy & multiple case study: 6 companies, 17 AI implementation cases, semi-structured interviews & case study & case observation (AI applications in SCOR processes) \\
\addlinespace[2pt]
S20 & Pham et al. (2024) & Journal of Business Research (3) & USA & 941 AI hospital-year obs + 941 matched non-AI hospital-year obs (246 matched hospitals), 40 US states, 2000-2020 & panel econometrics & adoption dummy (hospital AI adoption) \\
\addlinespace[2pt]
S21 & Sullivan \& Fosso (2024) & Journal of Business Research (3) & USA & two-wave survey, 107 US executives completing both waves (of 225 wave-1; 47.5\% effective RR), panel-recruited IT+business executives & survey-SEM & survey constructs (AI-enabled automation, AI-enabled analytics, AI-enabled relational capabilities) \\
\addlinespace[2pt]
S24 & Alam et al. (2025) & International Journal of Hospitality Management (3) & Malaysia & 336 respondents, stratified purposive sampling, hospitality industry & survey-SEM & survey construct (AI adoption + technology readiness) \\
\addlinespace[2pt]
S26 & Banna \& Alam (2025) & International Journal of Entrepreneurial Behaviour and Research (3) & EU & unbalanced panel of 1,479 firms (\textasciitilde{}12,900 obs), 26 EU countries, 2012-2023; CGM multi-way clustering + 2SLS-IV & panel econometrics & AI-related venture-capital funding received (log), OECD AI Policy Observatory data (turning point USD 11.3M); robustness: AI job intensity, AI patents, industry-level IV \\
\addlinespace[2pt]
S27 & Basnet et al. (2025) & International Review of Financial Analysis (3) & USA & US firms, 10-K filings 2005-2018 (pre-hype window); AI mentions classified actionable/speculative/irrelevant; causal design & panel econometrics & 10-K text (actionable vs speculative vs irrelevant AI disclosures) \\
\addlinespace[2pt]
S29 & Cao et al. (2025) & European Management Review (3) & USA & Study 1: instrument from 6,519 AI documents of 37 retailers; Study 2: fsQCA on survey of 140 U.S.-based retail managers & fsQCA & validated multi-level instrument (AI integration at task/function/firm level) \\
\addlinespace[2pt]
S32 & D’Amico et al. (2025) & Journal of Technology Transfer (3) & UK & 14,143 UK firms, 2004-2020 & panel econometrics & archival: Beauhurst-identified AI technology/product use, matched to Business Structure Database + UK Innovation Survey (2004-2020, biennial) \\
\addlinespace[2pt]
S34 & Fontanelli et al. (2025) & Journal of Economic Behavior and Organization (3) & France & French ICT survey 2019, firm-level; balancing/matching robustness & panel econometrics & survey (predictive AI use; buyers vs in-house developers) \\
\addlinespace[2pt]
S01 & Wamba-Taguimdje et al. (2020) & Business Process Management Journal (2) & multi-country & \textasciitilde{}500 vendor-published AI project mini-cases (IBM, AWS, Cloudera, Nvidia etc.), archival qualitative analysis & case study & case observation (AI-based transformation projects from vendor websites) \\
\addlinespace[2pt]
S10 & Sun et al. (2022) & Systems Research and Behavioral Science (2) & China & 234 AI-related manufacturers (GM/CEO respondents), CB-SEM & survey-SEM & survey construct (value co-creation in AI innovation ecosystem) \\
\addlinespace[2pt]
S11 & Ali et al. (2023) & Journal of Organizational Change Management (2) & UAE & single case (CMC Dubai): 9 semi-structured interviews with robotic-surgery team + archival data, triangulated & case study & case observation (adoption of AI/robotic surgery) \\
\addlinespace[2pt]
\end{longtable}

\clearpage
\begin{longtable}{@{}l p{2.4cm} l p{4.2cm} l p{6.2cm} p{6.2cm}@{}}
\caption{Included studies: outcomes, effect directions, and conditions (excerpt: 28 of 67 studies --- the 25 in the highest-ranked journals (AJG 4*, 4, 3) plus the first three AJG 2 studies; within a rating by year). Source: author's compilation.}\label{tab:coding-a2}\\
\toprule
ID & Study & Outcome & Outcome measure(s) & Direction & Conditions & Key finding \\
\midrule
\endfirsthead
\multicolumn{7}{@{}l}{\tablename~\thetable{} (continued)}\\
\toprule
ID & Study & Outcome & Outcome measure(s) & Direction & Conditions & Key finding \\
\midrule
\endhead
\bottomrule
\endfoot
S09 & Mishra et al. (2022) & Perf. & Perf: gross and net operating efficiency, net profitability, adspend, employment & mixed &  & AI focus reduces gross operating efficiency but increases net operating efficiency and profitability - automation raises costs more slowly than revenues. \\
\addlinespace[2pt]
S14 & Krakowski et al. (2023) & Both & Perf: competitive performance (game outcomes/ratings); CA: sources of persistent performance heterogeneity (= sustained advantage) across settings & conditional & substitution: traditional human capabilities become obsolete; complementation: new human-machine capabilities emerge that are unrelated or negatively related to traditional ones & AI adoption shifts the sources of competitive advantage - a new decision-making resource emerges at the human-AI intersection, unrelated to traditional capability. \\
\addlinespace[2pt]
S17 & Babina et al. (2024) & Perf. & Perf: growth in sales, employment, and market valuation & positive & channel: product innovation rather than process efficiency; gains concentrate among firms that were already large, which raises industry concentration (superstar dynamics) & AI-investing firms grow in sales, employment and value via product innovation; benefits concentrate among large firms, increasing concentration. \\
\addlinespace[2pt]
S07 & Leoni et al. (2022) & Perf. & Perf: perceived manufacturing firm performance (survey scale) & conditional & full mediation through knowledge management processes: the direct path from AI to manufacturing performance is not significant while both legs of the indirect path are; AI has no direct effect on supply chain resilience either; larger firms reach higher AI maturity & AI alone does not improve manufacturing firm performance - the effect exists exclusively through knowledge management processes (full mediation). \\
\addlinespace[2pt]
S45 & Adams et al. (2026) & Perf. & Perf: growth in sales, employment, assets, markups; stock-return response to monetary policy & positive & adoption concentrated in larger, more productive firms; adopters more responsive to monetary policy surprises & Firms adopting AI-powered pricing grow faster in sales, employment and markups - AI pricing is a capability large productive firms exploit first. \\
\addlinespace[2pt]
S46 & Agag et al. (2026) & Perf. & Perf: firm value; wage inequality as mediating channel & positive & wage inequality partially mediates: AI human capital compresses inequality and thereby raises firm value; the effect is stronger in dynamic and complex industries; by sector, the value effect is strongest in travel and the inequality compression in hospitality & AI human capital raises firm value partly by compressing wage inequality - a workforce-centered pathway to AI-enabled advantage, strongest in dynamic sectors. \\
\addlinespace[2pt]
S52 & Filieri et al. (2026) & Perf. & Perf: supply chain efficiency, sales performance, market value & conditional & inverted U: AI performance benefits diminish at extreme tariff exposure ("optimal investment threshold"); customer engagement needs strengthen the effect (boundary condition); tariff exposure induces AI integration & Trade tensions push firms into AI integration that improves efficiency, sales and value - up to a threshold where adversity overwhelms the compensation. \\
\addlinespace[2pt]
S02 & Baabdullah et al. (2021) & Perf. & Perf: perceived SME performance (survey scale) & positive & technology roadmapping (enabler, sig.); attitude (sig.); infrastructure readiness (sig.); awareness (sig.); professional expertise and technicality not significant & Acceptance of AI practices improves perceived SME performance; adoption is driven by roadmapping, attitude, infrastructure and awareness - not by professional expertise. \\
\addlinespace[2pt]
S03 & Chatterjee et al. (2021) & Both & Perf: perceived organizational performance (survey scale); CA: competitive advantage scale (COA1-3, Rogers-based; incl. market-share gain), discriminant-valid vs OP; predicted by performance & positive & mediators: B2B engagement, employee experience, information processing capability (all sig.); moderator: leadership support (sig., MGA); firm size/age/industry only controls (not moderators) & AI-CRM implementation improves perceived organizational performance and competitive advantage in B2B relationship management. \\
\addlinespace[2pt]
S04 & Chatterjee et al. (2022) & Perf. & Perf: firm performance scale (market+financial criteria, ROI/ROA, sales growth, market share, profitability) + B2B relationship satisfaction & positive & technology turbulence weakens the paths from automated decision-making and operational efficiency to B2B relationship satisfaction (negative moderator); leadership support strengthens the path from satisfaction to firm performance (positive moderator); the main effect of AI-CRM on firm performance is unconditional & AI-embedded CRM improves relationship satisfaction and firm performance; gains shrink under technology turbulence and grow with leadership support. \\
\addlinespace[2pt]
S05 & Hossain et al. (2022) & CA & CA: sustained competitive advantage = relative performance vs competitors over last 3 years (market share growth, sales growth, profitability, ROI) & conditional & mediators: market sensing, seizing and reconfiguring carry the effect of marketing analytics capability on sustained competitive advantage in full; moderator: AI adoption strengthens all three capability paths & Marketing analytics capability drives sensing/seizing/reconfiguring toward sustained competitive advantage; AI adoption amplifies this only when built on an analytics foundation. \\
\addlinespace[2pt]
S06 & Lee et al. (2022) & Perf. & Perf: revenue growth & conditional & AI-intensity threshold (low-level adoption: no effect); complementary technology investments (cloud, database) amplify; internal/exclusive R\&D strategy amplifies & AI pays off only above an adoption-intensity threshold, and more so with complementary technology investments and venture-specific internal R\&D. \\
\addlinespace[2pt]
S08 & Lui et al. (2022) & Perf. & Perf: firm market value (CAR; -1.77\% on event day) & negative & nonmanufacturing firms (more negative); weak IT capability (more negative); low credit rating (more negative) & investors react negatively to AI investment announcements on average; the penalty concentrates in nonmanufacturing firms with weak IT capability and low credit ratings. \\
\addlinespace[2pt]
S12 & Czarnitzki et al. (2023) & Perf. & Perf: firm productivity (labor- and value-added-based) & positive &  & AI adoption is robustly associated with higher firm productivity in representative German data; the effect grows with usage intensity. \\
\addlinespace[2pt]
S15 & Samadhiya et al. (2023) & Both & Perf: perceived firm performance (survey scale); CA: perceived sustainable competitiveness (survey scale) & positive & ICT capability (prerequisite); technological readiness (prerequisite); firm fit & AI-based partner relationship management improves firm performance and sustainable competitiveness, provided ICT capability and technological readiness are in place. \\
\addlinespace[2pt]
S16 & Wu \& Monfort (2023) & Perf. & Perf: perceived firm performance (survey scale) & positive & fsQCA configurations: marketing capabilities + customer value co-creation + market orientation + AI strategy jointly form necessary and sufficient recipes for high performance & AI marketing strategy raises firm performance, but only as part of configurations with marketing capabilities, co-creation and market orientation - a configurational result. \\
\addlinespace[2pt]
S18 & Cannas et al. (2024) & Perf. & Perf: qualitative competitiveness outcomes: costs, lead times, service level, quality, safety, sustainability & positive & barriers as boundary conditions: data quality, AI skills, high investment needs, unclear economic benefits, lack of AI cost-analysis experience & AI improves supply-chain competitiveness across SCOR processes, but realization depends on overcoming data-quality, skill and investment-evaluation barriers. \\
\addlinespace[2pt]
S20 & Pham et al. (2024) & Perf. & Perf: outpatient revenue, inpatient revenue, productivity, occupancy & positive & market share (larger-share hospitals adopt AI and benefit); endogeneity control essential - naive estimates misleading & Hospitals with larger market share adopt AI and gain in revenue, productivity and occupancy; without endogeneity controls the performance effects are misestimated. \\
\addlinespace[2pt]
S21 & Sullivan \& Fosso (2024) & Perf. & Perf: firm performance + process innovation + product innovation (survey scales) & conditional & adaptive response to market change is a full mediator, with no direct AI-to-performance paths modelled; environmental hostility weakens the automation path (negative moderator) and turns the relational path insignificant; environmental dynamism moderates as well; only ten of eighteen conditional indirect effects are significant & AI capabilities pay off through the firms adaptive response to market changes; environment hostility/dynamism condition how much each AI capability contributes. \\
\addlinespace[2pt]
S24 & Alam et al. (2025) & Both & Perf: value creation (survey scale); CA: sustainable competitive advantage (survey scale, modeled as mediator) & positive & sustainable competitive advantage mediates between AI adoption and value creation; dynamic capabilities first impede value creation through transition costs and become crucial later under turbulence; technological turbulence moderates only the path from capabilities to value creation & AI adoption creates value in hospitality mainly through sustainable competitive advantage as transmission channel; dynamic capabilities help only once transition costs are absorbed. \\
\addlinespace[2pt]
S26 & Banna \& Alam (2025) & Perf. & Perf: firm revenue growth & conditional & U-shaped effect with a quantified turning point at roughly USD 11.3M of AI venture funding: below it AI drags revenue, above it AI pays; coupling AI with an R\&D innovation strategy amplifies the effect, standalone AI is insufficient & AI investment follows a U-shaped payoff: short-term revenue drag, long-term gains - and only firms coupling AI with R\&D strategy capture the upside. \\
\addlinespace[2pt]
S27 & Basnet et al. (2025) & Perf. & Perf: firm value (market valuation); R\&D spending, patents, lagged productivity as channels & conditional & what counts is the substance of the disclosure: actionable disclosures bring valuation gains, innovation and lagged productivity, speculative or irrelevant ones bring nothing, and silent or vague peers are penalised & Markets reward only substantive (actionable) AI disclosures - early adopters with real implementation plans gain value via innovation; vague AI talk earns nothing. \\
\addlinespace[2pt]
S29 & Cao et al. (2025) & Perf. & Perf: efficiency and innovation outcomes (configurational) & conditional & configurational: no single AI function suffices - synergistic combinations (esp. customer service + cybersecurity with other functions); emergent capabilities: adaptive learning, predictive analytics, uncertainty mitigation & Retail performance comes from configurations of AI-enabled functions, not isolated AI uses - a directly configurational (fsQCA) result. \\
\addlinespace[2pt]
S32 & D’Amico et al. (2025) & Perf. & Perf: innovation performance (knowledge spillover of innovation) & conditional & complementarity with own R\&D investment (AI pays only if it complements internal R\&D); AI substitutes knowledge collaboration with certain external partners & AI adoption boosts innovation only when paired with the firms own R\&D investment - and partly replaces external knowledge collaboration. \\
\addlinespace[2pt]
S34 & Fontanelli et al. (2025) & Perf. & Perf: volatility of productivity growth (not the level) & conditional & AI buyers: higher volatility; in-house developers: none; complementary human capital (share of ICT engineers/technicians) mitigates the buyer effect & Off-the-shelf AI raises the riskiness of productivity outcomes; developing in-house or holding complementary ICT human capital neutralizes it. \\
\addlinespace[2pt]
S01 & Wamba-Taguimdje et al. (2020) & Perf. & Perf: organizational performance (financial, marketing, administrative) + process-level performance, qualitative & conditional & process reconfiguration required (mechanism: performance gains only when AI features are used to reconfigure processes); AI as configuration of complements (data, talent, domain knowledge, partnerships, infrastructure) & AI transformation projects improve organizational and process performance, but only when firms use AI to reconfigure their processes rather than overlay it on existing ones. \\
\addlinespace[2pt]
S10 & Sun et al. (2022) & CA & CA: perceived competitive advantage scale (own construct, Table 3) + innovation intelligibility; no performance construct in model & conditional & full mediation: value co-creation has no direct effect on competitive advantage; dynamic capabilities and innovation capabilities carry the entire effect & Value co-creation in AI ecosystems has no direct effect on competitive advantage - effects run entirely through dynamic and innovation capabilities. \\
\addlinespace[2pt]
S11 & Ali et al. (2023) & Both & Perf: four qualitative outcome categories: clinical (errors/infections down, patient experience), financial (quantified: \textasciitilde{}500K/month at 30 robotic cases), organizational (resource allocation), technological; CA: competitive position with explicit competitor comparison: patient choice over competitors, bargaining power with third-party payers, differentiation, upper hand over competition (P1) & positive & qualified robotic team incl. revenue-cycle team (capability prerequisite); cultural shift in medical community; regulatory approval (DHA, case-by-case); high acquisition costs not covered by insurance (entry barrier) & AI-enabled robotic surgery improves clinical, financial and technological outcomes and strengthens the competitive position of the healthcare organization. \\
\addlinespace[2pt]
\end{longtable}
\end{landscape}
\endgroup

\clearpage

% List of Aids deliberately NOT included in this draft — Arthur asks Fredrich
% first how AI usage should be documented (email of 29 July 2026).
% % List of Aids appendix — single source, input by main.tex AND the standalone
% supervisor drafts. Keep in sync with ai-usage-log.md (CLAUDE.md rule).
\section*{List of Aids}

\noindent The following AI tools and aids were used during the preparation of this thesis.
% Keep this table in sync with ai-usage-log.md — update immediately when a tool/section entry is added.

\vspace{0.5em}
\noindent\textbf{Global (all sections)}

\begin{longtable}{@{}p{2.9cm}p{3.2cm}p{9.0cm}@{}}
\toprule
\textbf{Section} & \textbf{Tool} & \textbf{Description} \\
\midrule
\endfirsthead
\toprule
\textbf{Section} & \textbf{Tool} & \textbf{Description} \\
\midrule
\endhead
\bottomrule
\endfoot
ALL & Grammarly Pro & For mistake and style checking throughout the thesis. Mistakes were corrected immediately; style suggestions were reviewed manually and applied when suitable. \\
\addlinespace
ALL & Claude Code (Anthropic) & Used throughout for LaTeX text correction and conversion, literature management (renaming PDFs, managing BibTeX), file structure, and drafting/editing section content. All AI-generated text was reviewed and approved before inclusion. \\
\addlinespace
ALL (writing phase) & library-rag --- local retrieval tool over the thesis corpus (built with Claude Code, 27 July 2026) & Local search system over the 67 full texts of the frozen corpus, used from the writing phase onward to locate citable passages. Runs entirely on the author's machine, no cloud services: Docling for PDF parsing, Qwen/Qwen3-Embedding-4B (fp8) for semantic search, BAAI/bge-reranker-v2-m3 for reranking, hybrid semantic and BM25 retrieval fused by reciprocal rank fusion. The tool retrieves and locates passages; it does not generate thesis text. Every verbatim quote it returns passes a two-stage programmatic check before it may enter a footnote: the quote must appear verbatim in the retrieved source passage, and its words must appear in the same order on the cited page of the PDF itself. The second stage also catches a quote attributed to the wrong passage or page. Its threshold was calibrated against 40 genuine and 60 fabricated quotes from this corpus (genuine: median 1.00; fabricated: maximum 0.18; the threshold of 0.60 admits none of the fabricated ones). Quotes the check cannot confirm are flagged for manual inspection of the page rather than passed through. Printed page numbers are calibrated per document from the printed page labels in the PDF and cross-checked against the Scopus page range; where no printed page number exists, the tool states this instead of supplying one. Corpus identification, screening, and coding were completed and frozen (17--18 July 2026) before this tool existed and are unaffected by it. \\
\addlinespace
Proposal & Claude Code (Anthropic) & Readability rewrite pass over the proposal prose; minimal simplifications merged into proposal\_final.tex after review; citations, quotes, and footnotes left unchanged. \\
\addlinespace
Proposal & Quillbot Paraphraser & Comparison paraphrase of citation-free proposal sentences (proposal\_quillbot.tex) to evaluate phrasing alternatives; citations, quotes, and footnotes left unchanged. \\
\addlinespace
Proposal & Claude Code (Anthropic) & Source check of the proposal against the collected article PDFs: every cited claim compared with the source text, page numbers and verification footnotes added to all citations, and five wordings corrected where they had drifted from the source. All corrections reviewed before merging. \\
\addlinespace
Proposal & Claude Code (Anthropic) & Shortening pass after I found the draft too dense: text condensed to 3 pages, one main statement per source, detail sentences removed. Citations and quotes kept verbatim; result reviewed by me. \\
\addlinespace
Method & Claude Code (Anthropic) & Corpus retrieval (17 July 2026): wrote and ran the script corpus/pull\_corpus.py, which executes the documented Scopus query via the Scopus API (WU license) and saves the frozen export (432 records, raw JSON + CSV). Verified that all three benchmark studies are retrieved and that record counts match Scopus. \\
\addlinespace
Method & Claude Code (Anthropic) & Title/abstract pre-screening (17 July 2026): Claude read the full abstracts of all 174 studies passing the AJG $\geq$ 2 filter and proposed include/exclude decisions with reason codes against the criteria catalogue approved beforehand; borderline cases flagged for full-text checking. All proposals recorded in corpus/screening\_2026-07-17.csv and reviewed by the author before becoming final. \\
\addlinespace
Method & Claude Code (Anthropic) & Full-text screening assistance (17--18 July 2026): Claude extracted and summarized method/measurement sections from the PDFs of all flagged borderline studies and proposed verdicts against the criteria catalogue; 12 studies excluded content-based, 1 confirmed as include; two inaccessible reports excluded as not retrieved (author's decision). Claude also matched, renamed, and de-duplicated the downloaded PDFs. All verdicts recorded with reasons and reviewed by the author. \\
\addlinespace
Results (coding) & Claude Code (Anthropic) + Gemini CLI (Google) & Dual independent AI coding of all 67 included studies against the frozen coding scheme (18 July 2026): Claude drafted all rows from the full texts (coder~1); Gemini 2.5 Pro coded the same studies blind---receiving only the coding scheme and PDF text, never Claude's drafts (coder~2). A documented self-consistency pass harmonized coder~1's effect-direction codes before comparison. Inter-coder agreement: method 90\%, outcome construct 88\%, effect direction 73\%. \\
\addlinespace
Results (coding) & Claude Code (Anthropic) & Adjudication support (18 July 2026): for all 29 coder disagreements and 10 random-sample agreement rows, Claude performed full-text reads and prepared evidence briefs quoting both coders' positions with page references; 15 factual gaps resolved from the full texts. The author adjudicated all disagreements against these briefs (18 individually, the remainder in a final batch review); decisions recorded per case, decision rules documented as written case law. No joint coder errors found in the sample checks; all 67 rows finalized after the author's batch approval. \\
\addlinespace
Theory (Section 2.1) & Claude Code (Anthropic) + library-rag & Drafting of Section 2.1 on the construct distinction (28 July 2026): citable passages located with the local retrieval tool; all four citations checked two-stage (verbatim in the source passage, word sequence on the cited PDF page) and their printed page numbers confirmed (pp. 215, 247, 550). Extended on 28 July 2026 by Barney's (1991, p. 102) definitions of competitive advantage and sustained competitive advantage, taken from the verified text sidecar and read off the page image; the corpus retrieval tool does not cover sources outside the corpus. The counts come from the fact sheet derived from the coding table, not from the tool. BibTeX entries generated from the frozen Scopus export with corpus/make\_bib.py. Human-voice pass applied. Text reviewed and edited by the author before inclusion. \\
\addlinespace
Results (descriptive mapping) & Claude Code (Anthropic) & Descriptive figures (18 July 2026): wrote and ran the script corpus/make\_figures.py, which derives five corpus-mapping charts (publication year, method, sample geography, outcome construct $\times$ effect direction, AJG rating) directly from the finalized coding table and writes them to figures/ as PDF. Category groupings are documented in the script; all counts verified to sum to $n = 67$. Colors validated as colorblind-safe. Figures reviewed by the author before inclusion. \\
\addlinespace
Theory (Sections 2.2--2.5) & Claude Code (Anthropic) + library-rag & Drafting of the four theory sections on general-purpose technology, the resource-based view, automation/augmentation and situated AI (28 July 2026). Citable passages from the corpus located with the local retrieval tool; all 11 corpus quotes checked two-stage (verbatim in the source passage, word sequence on the cited PDF page). Two quotes flagged by the check were inspected manually on the page; one revealed a page error (Krakowski et al., 1444 $\rightarrow$ 1445), which was corrected. The three theory anchors outside the corpus (Kemp 2024, Raisch \& Krakowski 2021, Brynjolfsson et al. 2021) are not in the retrieval index; their passages were read directly from the PDFs with pdfplumber and the printed page numbers verified against the running heads. Barney (1991) is cited from the verified text sidecar to the image scan (literature/barney1991firm.txt, OCR base checked against the scan page by page); both quoted passages were additionally read off the page images before the footnotes were set, and the original's printed typos are reproduced verbatim. BibTeX entries generated from the frozen Scopus export with corpus/make\_bib.py. Human-voice pass applied, including the paraphrase check against every footnote passage. Text reviewed and edited by the author before inclusion. \\
\addlinespace
Theory (Sections 2.1--2.5) & Claude Code (Anthropic) + Gemini CLI (Google) & Blind citation check of the finished theory chapter (28 July 2026), continuing the dual-AI principle from the coding phase. Claude wrote the script corpus/gemini\_verify\_citations.py, which parses every citation site in a section and splits the check by what can be decided mechanically: whether a quote appears verbatim on the cited page is settled deterministically by the local retrieval tool, the source-echo rule (three or more consecutive words shared between the prose and the cited passage) is computed by the script, and only the judgement of paraphrase fidelity goes to Gemini 2.5 Pro --- blind, receiving nothing but the sentence and the footnote passage. Of 25 citation sites in the chapter, the script flagged 7 for shared word sequences; the verdicts are recorded per site and collected in a report. The author decides every flag; the script proposes nothing and changes no text. \\
\addlinespace
Method / bibliography & Claude Code (Anthropic) & Citation rule for the seven studies without printed page numbers (28 July 2026): each of the seven checked against its publisher type and the rule approved by the author implemented --- article number in the bib eid field plus the article-internal page for the three Elsevier articles, section references for the four with publisher-relative or not-yet-assigned pagination. Two defects in bib.bib were found and fixed in the process: an unescaped ampersand in one title (would have aborted the LaTeX run at the first citation of it) and acronyms being lowercased by APA sentence casing, which had printed ``ai'', ``sme'' and ``b2b'' in the reference list. corpus/make\_bib.py now escapes ampersands and brace-protects acronyms, with self-tests; the 27 existing entries were corrected in one pass and the reference list checked in the compiled PDF. One Scopus metadata error corrected against the source PDF: the author of the European Journal of Marketing article was recorded as ``N., N. A.'' and is Tehrani, A. N. \\
\addlinespace
Results (descriptive mapping) & Claude Code (Anthropic) & Drafting of Section 4.1 (18 July 2026): Claude drafted the descriptive-mapping prose directly from the finalized coding table (aggregate counts only, no source citations) and applied the agreed human-voice pass (banned-pattern audit, voice matching). All numbers cross-checked against the coding table; text reviewed and edited by the author before inclusion. \\
\addlinespace
Appendix (coding table) & Claude Code (Anthropic) & Reader-facing cleanup of the appendix coding table (29 July 2026): a rendering layer in corpus/make\_appendix\_table.py strips coder shorthand from the printed table --- path coefficients, significance stars, hypothesis labels, model-fit values, and coder notes removed mechanically from 26 cells; 23 cells whose internal path abbreviations could not be resolved by rule were reworded into plain sentences, stored separately in corpus/appendix\_overrides.tsv. The coded values in corpus/coding\_table.csv are untouched --- the layer changes wording only, never a category, effect direction, or the identity of a condition. Two self-checks run at every regeneration: no shorthand token may survive into the rendered table, and every number printed must appear verbatim in the coded source cell. All 23 rewordings documented side-by-side (coded vs.\ appendix wording) for the author's review. \\
\addlinespace
Method (interim draft) & Claude Code (Anthropic) & First supervisor draft (29 July 2026): Claude assembled the standalone compile of the interim draft --- TOC skeleton, the full Method chapter, and the coding-table excerpt (the 25 studies in the highest-ranked journals plus three AJG 2 examples). Citation-verification footnotes are suppressed in the compiled draft (working files keep them); the drafted theory chapter and the List of Aids were left out of that round --- the former to keep the review load at one chapter, the latter pending the supervisor's guidance on how AI use should be documented. \\
\addlinespace
Method & Claude Code (Anthropic) & Drafting of the full Method chapter (29 July 2026): protocol prose assembled from the frozen method documentation (query wording, run date, AJG $\geq$ 2 rule, screening chain 432 $\rightarrow$ 174 $\rightarrow$ 81 $\rightarrow$ 67, E-code catalogue, dual-coding procedure with agreement statistics and the harmonization pass, pagination rule for the seven studies without printed page numbers). PRISMA-style flow diagram generated by corpus/make\_prisma.py from the screening-record counts, with checksum self-tests. One anchor citation (Rojon et al. 2021, p. 196) read column-wise from the PDF and footnoted per convention; a candidate passage on p. 195 was identified as a Tranfield quotation inside Rojon and rejected as a secondary source. The three benchmark studies are cited as entities without page numbers (no content claim). Human-voice pass applied. Text reviewed and edited by the author before sending. \\
\end{longtable}


\end{document}
